\section{Introduction}

\subsection{The Dirichlet divisor problem}

Let $d(n)$ denote the number of divisors of a positive integer $n$, and let
\begin{equation}
    D(x) := \sum_{n \leq x} d(n)
\end{equation}
be its summatory function. Dirichlet~\cite{Dirichlet1849} proved
in 1849 that
\begin{equation}
    D(x) = x \log x + (2\gamma - 1) x + \Delta(x),
\end{equation}
where $\gamma$ is the Euler--Mascheroni constant and
$\Delta(x) = O(x^{1/2})$. The Dirichlet divisor problem asks for the
infimum $\theta$ of all real numbers such that
$\Delta(x) = O(x^{\theta + \eps})$ for every $\eps > 0$.

Hardy~\cite{Hardy1916} proved that $\theta\geq 1/4$, and it is
conjectured that $\theta=1/4$. In the other direction,
Vorono\"i~\cite{Voronoi1903} proved the classical estimate
\begin{equation}
    \Delta(x)=O(x^{1/3}\log x),
\end{equation}
which implies $\theta\leq 1/3$. Further refinements eventually led
to the classical published bound of Huxley~\cite{Huxley2003}:
\begin{equation}
    \theta \leq 131/416.
\end{equation}

In this paper, rather than pursuing further improvements of $\theta$,
we explore an alternative approach: an exact combinatorial decomposition
of $D(x)$ via inclusion--exclusion on the number of distinct prime
factors of $n$.

\subsection{A new decomposition}

For a positive integer $n$, let $\omega(n)$ denote the number of distinct prime factors of $n$. For each integer $j\geq2$, define
\begin{equation}\label{eq:Tj-def}
    T_j(x) := \sum_{n \leq x}d(n)\binom{\omega(n)}{j}
\end{equation}
where $\binom{\omega(n)}{j}$ is the binomial coefficient. The main identity of this paper is the following.

\begin{theorem}\label{thm:identity}
    For every real number $x\geq1$,
    \begin{equation}\label{eq:main-identity}
        D(x)=1+\sum_{i\geq 1}(i+1)\pi(x^{1/i})+\sum_{j\geq 2}(-1)^j(j-1)T_j(x)
\end{equation}
where $\pi(x)$ denotes the prime-counting function. Both sums are
effectively finite for each fixed $x$: we have $\pi(x^{1/i})=0$
for all sufficiently large $i$, and $T_j(x)=0$ once $j$ exceeds
the maximum possible value of $\omega(n)$ among integers $n\leq x$.
\end{theorem}

The three terms on the right-hand side of \eqref{eq:main-identity} correspond, respectively, to $n = 1$, prime powers $\omega(n) = 1$, and integers with $\omega(n)\geq 2$. The proof of Theorem~\ref{thm:identity} is given in Section~\ref{sec:identity}.

Classical results of Dirichlet give the asymptotic
$D(x) \sim x \log x$. By the prime number
theorem~\cite{Apostol1976}, the prime-power contribution satisfies
\begin{equation}
     \sum_{i \geq 1} (i+1) \pi(x^{1/i}) \sim \frac{2x}{\log x},
\end{equation}
which is of strictly smaller order than $x \log x$. Hence the
first two terms in the identity are negligible compared with the
main asymptotic of $D(x)$, and the third sum
\begin{equation}
    \sum_{j \geq 2} (-1)^j (j-1) T_j(x)
\end{equation}
is responsible for the remaining main scale after cancellation.
Individually, each $T_j(x)$ grows rapidly with $j$, but the
alternating coefficients $(-1)^j (j-1)$ induce substantial
cancellation in the full finite sum. The fixed-$j$ asymptotics
proved below give information about individual terms, but they do
not by themselves determine this whole alternating sum.

\begin{example}
    At $x=10000$, direct computation gives
    \begin{align*}
        T_2(10000) &= 297{,}078 & T_3(10000) &= 140{,}640 \\
        T_4(10000) &= 26{,}674  & T_5(10000) &= 1{,}216
    \end{align*}
    and $T_j(10000)=0$ for $j\geq6$. The prime-power contribution equals
$2{,}711$, and the identity \eqref{eq:main-identity} yields
    \begin{align*}
        1+2{,}711+(T_2-2T_3+3T_4-4T_5)=1+2{,}711+90{,}956=93{,}668=D(10000),
    \end{align*}
    confirming Theorem~\ref{thm:identity} numerically. Moreover, the ratio between
$T_j(x)$ and the leading term later proved in Theorem~\ref{thm:main},
$\frac{2^j}{j!}x\log x(\log\log x)^j$, converges to 1 only very slowly: at $x=10^9$, this ratio is 0.43 for
$j=2$ and 0.21 for $j=3$. This slow convergence motivates the refined
asymptotic expansion of $T_j(x)$ developed in Section~\ref{sec:extraction}.
\end{example}

\subsection{Main results}

To analyze $T_j(x)$, we consider the two-variable Dirichlet series
\begin{equation}
    F(s, y) := \sum_{n \geq 1} \frac{d(n)(1+y)^{\omega(n)}}{n^s},
\end{equation}
which encodes $T_j(x)$ as the coefficient of $y^j$ in the generating
relation $\sum_{j \geq 0} T_j(x) y^j$. Here, we extend the definition
of $T_j(x)$ to all $j \geq 0$ by the formula~\eqref{eq:Tj-def}, so
that $T_0(x) = D(x)$ and $T_1(x) = \sum_{n \leq x} d(n) \omega(n)$.
The series $F(s, y)$ admits the factorization
\begin{equation}
    F(s, y) = \zeta(s)^{2+2y} H(s, y),
\end{equation}
where $H(s, y)$ is holomorphic in $s$ on a neighborhood of
$\Real(s) \geq 1$ and uniformly bounded for $y$ in a small disk
around $0$. This factorization isolates the singularity of
$F(s, y)$ at $s = 1$, and it is the starting point of the
Selberg--Delange method~\textup{[]}.

Our main asymptotic result is the following.

\begin{theorem}\label{thm:main}
For each fixed integer $j \geq 2$, as $x \to \infty$,
\begin{equation}
\begin{aligned}
    T_j(x)
    &= \frac{2^j}{j!} \, x \log x \, (\log \log x)^j \\
    &\quad + \frac{2^{j-1}}{(j-1)!}
       \bigl(\partial_y H(1, 0) - 2(1 - \gamma)\bigr)
       x \log x \, (\log \log x)^{j-1} \\
    &\quad + O\bigl(x \log x \, (\log \log x)^{j-2}\bigr).
\end{aligned}
\end{equation}
In particular,
\begin{equation}
    T_j(x) \sim \frac{2^j}{j!} \, x \log x \, (\log \log x)^j.
\end{equation}
\end{theorem}

The proof of Theorem~\ref{thm:main} is given in Sections~\ref{sec:perron}--\ref{sec:extraction}. The essential ingredients
are the truncated Perron formula, contour deformation around
$s = 1$, and a local Hankel analysis in the $w$-plane obtained
via the substitution $w = (s - 1) \log x$. The coefficient $\partial_y H(1, 0)$ admits an explicit expression as an absolutely convergent series over primes; we record this
expression in Section~\ref{sec:extraction}, where it is used in the proof of Theorem~\ref{thm:main}.

Theorem~\ref{thm:main} is a fixed-$j$ result. Combining it with
the identity~\eqref{eq:main-identity} does not by itself give an
asymptotic formula for $D(x)$, because the identity involves values
of $j$ that grow with $x$. In terms of the generating function
$A(x,y)=\sum_{j\geq0}T_j(x)y^j$, the alternating combination in
the identity is governed by values near $y=-1$, whereas the analytic
work below is carried out in a small disk around $y=0$. Recovering
$D(x)$ from the identity is therefore left as a future problem,
requiring estimates uniform in $j$ or a direct analysis near
$y=-1$.

\subsection{Outline of the paper}

The paper is organized as follows. In Section~\ref{sec:identity},
we prove the identity (Theorem~\ref{thm:identity}) by partitioning
$D(x)$ according to $\omega(n)$ and applying an inclusion--exclusion
argument to the composite part. In Section~\ref{sec:dirichlet-series},
we introduce the two-variable Dirichlet series $F(s, y)$ and establish
its factorization $F(s, y) = \zeta(s)^{2+2y} H(s, y)$, together with
the analytic properties of $H(s, y)$ needed in the sequel.

Sections~\ref{sec:perron}--\ref{sec:extraction} are devoted to the
proof of Theorem~\ref{thm:main}. In Section~\ref{sec:perron}, we
apply the truncated Perron formula to the generating partial sum
$A(x, y) := \sum_{n \leq x} d(n) (1+y)^{\omega(n)}$ and deform the
contour around $s = 1$. In Section~\ref{sec:hankel}, we carry out
the local Hankel analysis via the substitution $w = (s-1) \log x$
and evaluate the main and lower-order contributions. In
Section~\ref{sec:extraction}, we extract the asymptotic formula for
$T_j(x)$ from that of $A(x, y)$ by differentiating in $y$, thereby
completing the proof of Theorem~\ref{thm:main}.

Finally, in Section~\ref{sec:outlook}, we discuss why the fixed-$j$
asymptotics and the identity do not yet determine the full
alternating sum for $D(x)$, and what additional uniform estimates
would be needed.

% ===== Section 2 =====
